\documentclass[aspectratio=169, 11pt]{beamer}
\usepackage{beamer_preamble}

\begin{document}

% ==== Title ===============================================================
\titleslide{Lesson 4-1 \textperiodcentered\ Period 1}{Inverse Variation --- Introduction}

% ==== Do Now ==============================================================
\begin{frame}{Do Now --- Notice \& Wonder}{Algebra 2 \textperiodcentered\ Lesson 4-1 \textperiodcentered\ Period 1}
\phasebadges{1}{5}{notice \& wonder}
\vspace{0.1cm}
\begin{projcard}{Two rectangles --- both have area = 144 sq units}{slidegold}
\small
Rect 1: length = 72, width = 2 \hspace{2em} Rect 2: length = 24, width = 6

\vspace{0.15cm}
\begin{enumerate}
  \setlength{\itemsep}{1pt}
  \item \textbf{Notice:} Write something you notice.
  \item \textbf{Wonder:} Write something you wonder about rectangles with area 144.
  \item \textbf{Sketch:} Draw a third rectangle with area 144. Label length and width.
\end{enumerate}
\end{projcard}
\end{frame}

% ==== Launch ==============================================================
\begin{frame}{Launch --- Inverse Variation Rules}{Algebra 2 \textperiodcentered\ Lesson 4-1 \textperiodcentered\ Period 1}
\phasebadges{1-2}{12}{teacher models}
\vspace{0.1cm}
\begin{projcard}{Inverse Variation Rules (keep on your page)}{slidegreen}
\small
\begin{enumerate}
  \setlength{\itemsep}{1pt}
  \item INVERSELY: their \textbf{product} is a constant \(k\) \hspace{1em} \(xy = k\) \hspace{0.5em} or \hspace{0.5em} \(y = k/x\)
  \item As one variable INCREASES, the other DECREASES proportionally.
  \item To find \(k\): multiply any given pair.
  \item To find unknown: write \(y = k/x\), then substitute.
\end{enumerate}
\vspace{0.1cm}
\textcolor{slideaccent}{\textbf{TRAP:}} \(y/x = k\) = DIRECT \hspace{2em} \(xy = k\) = INVERSE
\end{projcard}
\end{frame}

% ==== Launch Example 1 ====================================================
\begin{frame}{Launch --- Example 1: Table Check}{Algebra 2 \textperiodcentered\ Lesson 4-1 \textperiodcentered\ Period 1}
\phasebadges{1-2}{12}{teacher models}
\vspace{0.2cm}
\begin{projcard}{Is this an inverse variation?}{slidegreen}
\begin{tabular}{l|cccccc}
\(x\) & 1 & 2 & 3 & 5 & 6 & 15 \\
\(y\) & 25.5 & 12.75 & 8.50 & 5.10 & 4.25 & 1.70 \\
\end{tabular}

\vspace{0.3cm}
\textbf{Step:} Compute \(xy\) for each column.\\
If every product equals the same number \(\rightarrow\) \textbf{YES, inverse variation.}\\[0.15cm]
\textbf{Bridge:} What was \(xy\) for your Do Now rectangles?
\end{projcard}
\end{frame}

% ==== Launch Example 2 ====================================================
\begin{frame}{Launch --- Example 2: Write the Equation}{Algebra 2 \textperiodcentered\ Lesson 4-1 \textperiodcentered\ Period 1}
\phasebadges{1-2}{12}{teacher models}
\vspace{0.1cm}
\begin{projcard}{Given a pair of values, write the equation}{slidegreen}
\small Two variables vary inversely. \(x = 10\) when \(y = 3\).

\vspace{0.15cm}
\begin{enumerate}
  \setlength{\itemsep}{1pt}
  \item Find \(k\): \(k = xy = \) \underline{\hspace{2cm}}
  \item Write the equation: \(y = k/x\ \longrightarrow\ y = \) \underline{\hspace{2cm}}
  \item Find \(y\) when \(x = -6\): \(y = \) \underline{\hspace{2cm}}
\end{enumerate}
\vspace{0.1cm}
\small\textcolor{slidegray}{30-sec partner check: What is \(k\) for your Do Now rectangles (area = 144)?}
\end{projcard}
\end{frame}

% ==== Explore =============================================================
\begin{frame}{Explore --- TPS Practice}{Algebra 2 \textperiodcentered\ Lesson 4-1 \textperiodcentered\ Period 1}
\phasebadges{1-2}{33}{student work}
\small\textcolor{slidegray}{7 items in order. \(\sim\)4--5 min per item. Compute \(xy\) for every table before answering.}

\vspace{0.1cm}
\begin{projcard}{Try It 1 \textperiodcentered\ Two tables: inverse variation?}{slideblue}
Compute \(xy\) for each column. Is the product constant?
\end{projcard}
\vspace{-0.2cm}
\begin{projcard}{Practice \#9 \textperiodcentered\ Generalize}{slideblue}
From the table, how can you tell the data do NOT show inverse variation?
\end{projcard}
\vspace{-0.2cm}
\begin{projcard}{Practice \#14 \textperiodcentered\ Table check}{slideblue}
Do these values represent inverse variation? Explain.
\end{projcard}
\vspace{-0.2cm}
\begin{projcard}{Try It 2 \textperiodcentered\ Write the equation}{slideblue}
\(x = 6\) when \(y = 1/2\). (a) Write the equation. (b) Find \(y\) when \(x = 15\).
\end{projcard}
\vspace{-0.2cm}
\begin{projcard}{Practice \#13 / \#16 / \#10 \textperiodcentered\ Generalize \& Apply}{slideblue}
Write \(k\) in terms of \(x\) and \(y\). Apply. Extension: why can \(x \neq 0\)?
\end{projcard}
\end{frame}

% ==== Share / Summary =====================================================
\begin{frame}{Share / Summary}{Algebra 2 \textperiodcentered\ Lesson 4-1 \textperiodcentered\ Period 1}
\phasebadges{2}{5}{academic conversation}
\vspace{0.2cm}
\begin{projcard}{Sentence frame \textperiodcentered\ Pair share then whole class}{slideblue}
``As \underline{\hspace{1.2cm}} increases, \underline{\hspace{1.2cm}} decreases.
Their product is \underline{\hspace{1.2cm}}, so the equation is
\(y = \underline{\hspace{1cm}}/x\).''

\vspace{0.2cm}
Class signals \textbf{thumbs up} (agree) / \textbf{sideways} (unsure).
\end{projcard}

\vspace{0.15cm}
\begin{projcard}{Bridge to Period 2}{slideaccent}
Next class: \textbf{real-world inverse variation problems} (bouzouki string).\\
Also: DIRECT vs INVERSE --- both can appear in the same context!\\
Period 2 carries the DOK-3 performance task.
\end{projcard}
\end{frame}

% ==== Exit Ticket =========================================================
\begin{frame}{Exit Ticket --- Summary Recap}{Algebra 2 \textperiodcentered\ Lesson 4-1 \textperiodcentered\ Period 1}
\phasebadges{1-2}{2}{not graded}
\vspace{0.2cm}
\begin{projcard}{Complete on your exit ticket}{slidegold}
\begin{enumerate}
  \setlength{\itemsep}{0.2cm}
  \item Today I learned that \underline{\hspace{5cm}}.
  \item The difference between inverse and direct variation is \underline{\hspace{4cm}}.
  \item One thing I'm still unsure about is \underline{\hspace{5cm}}.
\end{enumerate}
\end{projcard}
\end{frame}

\end{document}
